\chapter{Testing}
Up until now, diverse methods of integrating an OpenArena server have been detailed, but the testing part is yet to be tackled.
It has a significant role, ensuring the correctness of the implementations.
\section{Introducing OpenArena Client}
Given that OpenArena has a client-server design, it is important to briefly explain the client side too, for a rounded understanding of the infrastructure. \\
The OpenArena client is primary represented by the audio and video graphics rendering libraries. Its core library is OpenGL\footnote{\url{https://www.opengl.org/}}, but
the used version is an old one. \\
As the server runs in an virtual environment provided by KVM/QEMU and the clients run locally, there is some networking configuration done. Therefore, the server side
uses bridging to enable the guest system to appear on the same network as the host. This is achieved by setting up a bridge interface and attaching the unikernel
network interface to it, similar to what the figure\footnote{\url{https://www.techtarget.com/searchitoperations/tip/Adding-an-additional-virtual-bridge-to-a-KVM-environment}} below suggests.
\fig[scale=0.75]{./src/img/bridging.pdf}{img:bridging}{Role of Virtual Bridge}
\section{Gameplay Simulation}
A solution for testing the server is interacting with the clients, so called game players.
In order to achieve it, a Python script has been used to simulate a gameplay environment, which leverages the Pexpect module.
That is, Pexpect spawns child application and controls them as if a human were typing commands. \\ \\
Therefore, the script creates a thread for each client and employs Pexpect to generate the OpenArena clients.
A client stays connected for 30 seconds. During this time, it does certain actions: moves in a direction, shoots for a short amount of time.
After this, a delay time of 15 seconds takes place and then the cycle start again.
\fig[scale=1]{./src/img/testing.drawio.pdf}{img:testing}{Testing Environment}